%=============================================================
\chapter{Methodology}
\label{ch:method}
%=============================================================

% ─── 章節破題 ───────────────────────────────────────────────
本章說明如何設計實驗以回答第~\ref{ch:motivation}~章提出的四個研究問題。
一個好的實驗設計必須同時滿足三個條件：
可重現性（所有關鍵參數均有明確記錄）、
公平性（對照系統在同等資源限制下比較）、
以及可解釋性（每個實驗的結論能夠獨立歸因，而非多個因素同時變動）。
PaceANN 的評估涉及三個在架構上有本質差異的系統，
公平性問題尤為關鍵，本章以獨立節次（§\ref{sec:fairness}）專門討論。

%-------------------------------------------------------------
\section{實驗假設與 RQ 對應}
\label{sec:hypotheses}
%-------------------------------------------------------------

每個子實驗對應第~\ref{ch:motivation}~章的一個研究問題，
並有明確的可驗證假設：

\begin{description}
  \item[\textbf{E1}（對應 RQ1）]
    \emph{PQ 距離比值能否可靠偵測 beam search 的收斂？}
    假設：在相同 recall 目標下，啟用 ET（$\theta > 0$）的平均 SSD I/O 數顯著低於停用 ET；
    且 recall 損失隨 $\theta$ 單調遞增。

  \item[\textbf{E2}（對應 RQ2）]
    \emph{快取覆蓋如何影響 ET 的觸發可靠性？}
    假設：快取覆蓋較充分時，相同 $\theta$ 的 recall 損失較小，
    亦即快取是 ET 可靠性的前提。

  \item[\textbf{E3}（對應 RQ3）]
    \emph{在等記憶體預算下，PaceANN 的 QPS-recall Pareto 是否優於 Starling？}
    假設：在相同的記憶體預算內，
    PaceANN 在 recall $\geq 0.95$ 的區間達到更高 QPS，
    且 P99 延遲不劣於 Starling。

  \item[\textbf{E4}（對應 RQ4）]
    \emph{$\theta$ 的最優值在不同 recall 目標下如何變化？}
    假設：存在一個資料集相關的最優 $\theta^*$，
    使 QPS-recall trade-off 最佳化；
    $\theta^*$ 與 PQ 壓縮誤差分布相關，可透過掃描估計。
\end{description}

%-------------------------------------------------------------
\section{實驗環境}
\label{sec:testbed}
%-------------------------------------------------------------

\begin{table}[h]
\centering
\caption{實驗硬體規格}
\label{tab:hardware}
\begin{tabular}{ll}
\toprule
元件 & 規格 \\
\midrule
CPU & AMD Ryzen 7 9700X（8 核心 / 16 執行緒，Zen 5） \\
DRAM & 64\,GB DDR5 \\
NVMe SSD & Predator GM9000 1\,TB（PCIe Gen 5） \\
作業系統 & Linux 7.0 \\
非同步 I/O & \texttt{libaio}（\texttt{O\_DIRECT | O\_RDONLY}） \\
\bottomrule
\end{tabular}
\end{table}

所有實驗均在同一台機器上進行，避免跨機器的硬體差異影響可重現性。
Worker thread 數固定為 $T = 16$（與 CPU 邏輯核心數一致，避免過度訂購）。
SSD 在實驗前未執行其他 I/O 密集工作，確保快取狀態一致。

%-------------------------------------------------------------
\section{資料集與索引設定}
\label{sec:dataset}
%-------------------------------------------------------------

主要實驗使用 100M 級的磁碟式資料集（例如 SIFT100M：1 億筆 128 維 uint8 向量，原始大小約 12.8\,GB）。
選這個規模，是因為它夠大、在有限記憶體預算下 SSD 儲存無可避免，正好對應本研究關注的磁碟式情境；
uint8 型別也與生產環境的影像特徵一致（注意 SIFT1M 為 float32，不可混用），而且是社群常用的 benchmark，方便和文獻交叉比對。
查詢集使用官方的 10,000 筆 query、以 top-10 精確最近鄰為 ground truth；所有量測都用完整 10,000 筆、不做截斷，以確保 recall 估計的統計穩定。

DiskANN 與 PaceANN 的 Vamana 索引以表~\ref{tab:index-build} 的參數建構。

\begin{table}[h]
\centering
\caption{DiskANN / PaceANN Vamana 索引建構參數}
\label{tab:index-build}
\begin{tabular}{lll}
\toprule
參數 & 數值 & 說明 \\
\midrule
最大出度 $R$ & 96 & Vamana Robust Pruning 的圖密度上限 \\
建構搜尋列表 $L_\text{build}$ & 120 & 建構時的 beam search 預算 \\
分塊大小 $B$ & 2 & 磁碟 sector 每塊包含的節點數 \\
合併輪次 $M$ & 40 & 增量建構的合併參數 \\
Medoid 數量 $K$ & 5 & 多入口點數（分片均等預算） \\
\bottomrule
\end{tabular}
\end{table}

$R=96$ 是在圖密度與索引大小之間取平衡：$R=128$ 雖然連接更豐富，但 sector 大小固定，$R$ 越大代表每次要讀的 sector 越多，反而可能拖低 I/O 效率。

Starling 使用相同的圖建構參數（$R=96$、$L_\text{build}=120$），但額外做了 block shuffling，把相鄰節點打包到同一磁碟頁面；
其 navigation graph（mem-index）以 $R_\text{mem}=20$、$L_\text{mem}=128$ 建構，涵蓋約 0.1\% 的隨機採樣節點。
要說明的是，Starling 的索引格式與 DiskANN 並不相容，兩者無法共用同一個索引檔，因此圖的建構細節上仍有些微差異；這是現有評估方法的固有限制，我們在呈現比較時都會標明。

%-------------------------------------------------------------
\section{公平性協定：記憶體預算正規化}
\label{sec:fairness}
%-------------------------------------------------------------

PaceANN 與 Starling 在架構上的根本差異在於：

\begin{itemize}
  \item DiskANN / PaceANN 的最小啟動記憶體僅需 PQ 壓縮向量（約 1.74\,GB），
    剩餘預算可彈性分配給 BNC。
  \item Starling 在啟動時強制載入 \texttt{partition.bin}（827\,MB）到兩個記憶體結構
    （\texttt{id2page\_} 約 400\,MB、\texttt{gp\_layout\_} 約 400\,MB），
    加上 navigation graph（約 31\,MB），
    最小 RSS 固定為 \textbf{3.07\,GB}，無法降低。
\end{itemize}

若以固定絕對預算（例如「同為 3\,GB」）比較，
Starling 恰好用完而 PaceANN 還有餘裕，對 PaceANN 有利；
若以不同絕對預算分別取最優，
則兩者資源不等，比較失去意義。

為了讓四個系統站在同一條起跑線上，本研究以固定的總記憶體預算（涵蓋 PQ 與快取）作為公平基準，各系統都把記憶體用到接近這個上限、但不超過。
這樣比較的意義，在於它反映的是「相同硬體成本下誰更快」，而不是讓某個系統偷偷多用記憶體去換效能。各系統在該預算下的實際 RSS 需在量測機器上校準後填入（見第~\ref{ch:evaluation}~章）。

量測時一律採冷啟動、不做暖機。server 全新啟動後，直接以原始查詢集完整跑一次，不重複、不循環、不預熱，再記錄結果。
這一點其實很關鍵：PaceANN 的動態快取若以重複查詢預熱，命中率會被人為灌高、I/O 趨近於零，量到的就不是真實流量下的效能。
因此本研究刻意避免任何 warmup，讓 PaceANN 的快取只透過 BFS 冷啟動種子與正式查詢過程中的 demand-fill 自然累積，這樣的數字才代表實際部署時的表現。

執行緒方面，所有系統都設 $T = 16$ 個 worker thread，beam width 依各系統的最優設定
（DiskANN 與 PaceANN 用 $W = 8$，Starling 用 $W = 4$，與其原論文一致）。

%-------------------------------------------------------------
\section{量測方法論}
\label{sec:measurement}
%-------------------------------------------------------------

延遲數字（平均、P50、P99、P999）都在 server 端計算，時間戳從 worker thread 完整收到查詢封包後開始、到送出結果前結束，
排除 TCP 建連與網路傳輸，讓延遲只反映搜尋演算法本身的成本。QPS 則在 client 端以牆鐘時間計算，從發出第一筆查詢到收到最後一筆回應，反映端到端的吞吐量。

主要結果以並發度 $C \in \{1, 8, 32\}$ 呈現：$C=1$ 是單客戶端的串列延遲基準，$C=8$ 是中等並發的典型 server 負載，$C=32$ 則是高並發壓力（$C>T$，存在查詢排隊）。
其中 $C=8$ 是主要的 QPS 比較點，因為此時吞吐量與尾延遲都落在有意義的工作區間。

搜尋長度 $L$ 是控制 recall 的主要旋鈕：$L$ 越大 recall 越高，但每筆查詢的 hop 數與 I/O 也越多。
藉由掃描 $L \in \{20, 30, 50, 80, 100, 150, 200, 250, 300\}$，可在 recall 軸上畫出 QPS--recall 的 Pareto 曲線，涵蓋 recall 約 0.60 到 0.99 的區間。

%-------------------------------------------------------------
\section{消融研究設計}
\label{sec:ablation}
%-------------------------------------------------------------

為將快取貢獻與 ET 貢獻分開量化，
實驗採用漸進式啟用策略：

\begin{table}[h]
\centering
\caption{消融研究的系統設定組合}
\label{tab:ablation}
\begin{tabular}{lll}
\toprule
設定代號 & 動態快取（BNC） & Early Termination \\
\midrule
\texttt{baseline}    & \texttimes & \texttimes \\
\texttt{cache\_only} & \checkmark & \texttimes \\
\texttt{ET\_only}    & \texttimes & \checkmark \\
\texttt{full}        & \checkmark & \checkmark \\
\bottomrule
\end{tabular}
\end{table}

每組設定的 QPS-recall Pareto 差距揭示對應元件的邊際貢獻：
\begin{itemize}
  \item \texttt{cache\_only} vs. \texttt{baseline}：
    量化動態快取命中率帶來的 I/O 節省。
  \item \texttt{ET\_only} vs. \texttt{baseline}（無快取）：
    在沒有快取保護的情況下，ET 的 recall 損失（反映假陽性終止的嚴重性）。
  \item \texttt{full} vs. \texttt{cache\_only}：
    在快取已覆蓋收斂區的前提下，ET 的額外貢獻。
  \item \texttt{full} vs. \texttt{ET\_only}：
    量化快取對 ET 可靠性的「使能效應」。
\end{itemize}

最後一個比較（\texttt{full} vs. \texttt{ET\_only}）是
§\ref{sec:codesign} 協同必要性論點的直接實驗驗證：
若協同效應真實存在，兩者的 recall 差距應在相同 $\theta$ 下顯著可見。

%-------------------------------------------------------------
\section{$\theta$ 參數敏感性研究}
\label{sec:theta-sweep}
%-------------------------------------------------------------

ET 的核心參數 $\theta$ 決定停止條件的保守程度。
本研究以以下方法論系統化地確定其最優設定範圍：

我們固定快取預算，掃描 $\theta \in \{1.05, 1.10, 1.12, 1.20, \infty\}$（$\infty$ 表示停用 ET），對每個 $\theta$ 跑完整的 $L$-sweep（$L \in \{20, 30, 50, \ldots, 300\}$），
記錄 QPS、Recall@10 與每筆查詢的平均 SSD I/O 數（MeanIOs）。MeanIOs 是關鍵的中間指標：它量化了 ET 實際擋掉多少 SSD 讀取，讓 QPS 的改善能歸因於「更少 I/O」，而不是其他因素。

為了避免用更大的 $L$ 去補償激進 $\theta$ 造成的 recall 損失而產生對比偏差，我們對每個 recall 門檻 $r^*$，把 Pareto 最優的 $\theta$ 定義為在滿足 recall $\ge r^*$ 的前提下使 QPS 最大者：

\begin{equation}
  \theta^*(r^*) = \argmax_{\theta} \text{QPS}(\theta, L^*(r^*,\theta))
  \quad \text{s.t. } \text{Recall}(\theta, L^*(r^*,\theta)) \geq r^*
  \label{eq:theta-opt}
\end{equation}

其中 $L^*(r^*, \theta)$ 是在給定 $\theta$ 下達到 recall $r^*$ 所需的最小 $L$。

理論上 $\theta$ 的容忍帶應該和 PQ 的量化誤差成正比：PQ 距離估計的偏差越大，$\theta$ 就該設得越保守；偏差小時 $\theta$ 可以接近 1。
我們藉由比較不同 $\theta$ 下 recall 的下降速率，反推本實驗設定下等效的 PQ 誤差容忍帶，作為未來其他資料集選 $\theta$ 的參考。

%-------------------------------------------------------------
\section{量測指標定義}
\label{sec:metrics}
%-------------------------------------------------------------

\begin{description}
  \item[Recall@10]
    $|S_{10}(q) \cap G_{10}(q)| \;/\; 10$，
    其中 $S_{10}(q)$ 為系統返回的 top-10，$G_{10}(q)$ 為 ground truth top-10。
    以 10,000 筆 query 的算術平均值報告。

  \item[QPS（Queries Per Second）]
    $n_\text{total} / T_\text{wall}$，
    以 client 端牆鐘時間計算，
    包含 TCP 傳輸延遲（通常 $< 0.1\,\text{ms}$，在 localhost 模式下可忽略）。

  \item[Server-side Latency（Mean / P50 / P99 / P999）]
    以 server 端的處理時間計算（不含 TCP 傳輸）。
    P99 與 P999 是評估尾延遲行為的主要指標，
    對應 SLA 中「99\% / 99.9\% 請求須在時限內完成」的要求。

  \item[MeanIOs]
    每次查詢觸發的 SSD 讀取次數的平均值（不含快取命中）。
    此指標僅適用於 PaceANN（DiskANN 系），
    用於量化 ET 的 I/O 節省效果。

  \item[RSS（Resident Set Size）]
    Server process 在量測期間的常駐記憶體大小，
    以 \texttt{/proc/\$PID/status} 的 \texttt{VmRSS} 欄位讀取。
    用於驗證實際記憶體用量是否在預算內。
\end{description}

%-------------------------------------------------------------
\section{可重現性說明}
\label{sec:reproducibility}
%-------------------------------------------------------------

為便於重現，所有 benchmark 腳本都以 shell script 保存，包含明確的 server 啟動命令、並發數設定與輸出 CSV 格式。
每次 server 啟動後會等待數十秒，確保 index 完整載入後才開始量測。

索引建構使用固定隨機種子，確保 Vamana 圖的可重現性。
Starling 的索引在獨立機器上建構後以 \texttt{rsync} 傳輸，
其 medoid 節點（25,384,254）與 DiskANN 索引（94,370,396）不同，
此差異已在分析中明確標記。
